\section{The Merits Firewall and Administrable Boundaries}
\label{sec:objections}

A front-door rule succeeds because its limiting principles are part of the rule. Ordinary law retains the merits. A closed, staged record binds information cost. Existing legal interests define who may invoke access. These allocations yield a sectoral information rule with a risk-specific control lane, protected access, and prompt termination for enterprises outside that lane.

This Part develops those affirmative boundaries and then tests them against fragmented systems, user and third-party intervention, arbitration, cross-border suppliers, and speech. The front door makes facts available; the governing law decides what those facts mean.

\subsection{Ordinary Law Receives Better Factual Predicates}

Courts have long handled aircraft, pharmaceuticals, industrial equipment, distributed supply chains, proprietary software, and corporate records. Pleading on information and belief, early discovery, party requests, subpoenas, expert testing, circumstantial proof, amendment, and sanctions can address severe information asymmetry. AI receives no bespoke merits regime simply because its technology is complex.

The proposal accepts that premise. It creates no general AI tort, standard of care, product classification, causation presumption, or damages rule. Abraham and Sharkey's account of objective reasonableness explains how ordinary negligence can evaluate many AI disputes without reconstructing every internal model state.\lrfootnote{See \artcite{abrahamSharkey2027}.} The cascade in Part I identifies observable deployment conduct precisely because the merits can remain ordinary.

The distinction lies in procedural prerequisites. Rule 8 and Rule 11 can accommodate incomplete evidence, but the complaint still requires a legal claim and enough facts to make it plausible. Rule 26 begins after a complaint has identified a litigable target. Rule 34 reaches a party's possession, custody, or control; Rule 45 reaches a known and served nonparty. Early discovery requires a recipient and judicial authorization. \emph{Strike 3} shows how well the path can work when a claimant has a unique incident identifier and a known intermediary with the identity record. A divided AI deployment can withhold one of those predicates. Doe pleading cannot serve an unknown entity, relation back depends on notice and mistake, and Rule 27 preserves known testimony. Those tools become powerful after actor, version, and record holder are visible.

The front door fills that anterior stage. The operator already knows which system it incorporated into a consequential decision or made available to the claimant. The ex ante record duty requires provenance and risk control to remain knowable. The first answer passes that bounded information into ordinary procedure. From then on, the Federal Rules and the substantive cause of action perform their familiar work.

The rule also supplies a function outside discovery: creation. The ex ante duty generates release approval and version history while the deployment is operating. Litigation preservation begins when a dispute becomes reasonably foreseeable and protects that earlier record. Sectoral recordkeeping therefore makes consequential activity observable before the incentive to reconstruct it appears.

That is a familiar regulatory move. Aviation safety, financial transactions, employment records, medical care, and public administration use ex ante records because later adjudication and oversight require them. The justification comes from the protected sector and decision, not from a claim that AI is uniquely inscrutable. A designated deployment receives a front door when the legislature concludes that its consequences and divided information warrant one.

Existing attribution doctrine remains fully available. A corporation can be directly liable for its own design, supervision, warning, monitoring, or deployment choices; agency law can attribute conduct; products law can operate without a culpable engineer; electronic-transactions law can give effect to authorized automated acts.\lrfootnote{See \bluecite[\S\S~7.03(1), 7.05, 7.07]{restatementAgency2006}; \bluecite[\S\S~1--2]{restatementProducts1998}; \bluecite[\S~14]{ueta1999}.} In \emph{Mobley}, the court used a statutory agency theory at pleading rather than waiting for a legislative AI fix. The front door is valuable because it allows claimants and defendants to identify the facts on which such ordinary theories succeed or fail.

\subsection{Staged Access Binds Information Cost}

A firm may prevail on duty, defect, causation, damages, immunity, or another defense after paying for record design, retention, internal investigation, privilege review, trade-secret protection, third-party notice, and access. Multiple enterprises can duplicate work. The rule therefore assigns expense where standardization lowers total cost and screens requests before deeper production.

Geistfeld's analysis of insurability and overlapping protection warns against liability rules that raise prices for the same people they aim to protect.\lrfootnote{See \artcite{geistfeldInsurability2026}.} The front door avoids the largest version of that problem by preserving the substantive class of compensable losses and ultimate burdens. Its remaining information cost should be defended directly: a limited record is the institutional price of making a designated high-impact deployment observable.

The price is controlled at six points. First, the legislature designates sectors and protected interests; capability or marketing alone never activates the regime. Second, the claimant identifies an existing claim or equitable entitlement and pleads concrete, risk-specific facts. Third, Stage One is limited to identity, versions, control allocation, incident events, preservation, and governing policies. Fourth, Stage Two follows when the initial record supports claim-specific need and is bounded by time, version, population, feature, and risk. Fifth, source code, weights, sensitive data, and security controls require necessity and consideration of less intrusive proof. Sixth, an enterprise without material control completes the first answer and terminates its special role.

Cost allocation can reinforce the screen. The operator bears the ordinary expense of maintaining and retrieving the minimum record because observability is a condition of covered deployment. A requester bears the cost of extraordinary formatting, duplicative demands, or an expert protocol when the tribunal finds the expense disproportionate to the defined interest. Expenses caused by an incomplete or evasive answer shift to the nonperforming enterprise. Illinois Rule 224 offers a useful identity-stage analogue by limiting the proceeding to identification and assigning reasonable compliance costs to the petitioner; a sectoral statute can modify that baseline where the operator's ex ante duty makes ordinary retrieval part of compliance.\lrfootnote{See \statcite{illinoisRule224}.}

The ex ante record can also lower total cost. A deployment ID, defined owners, version lineage, event schema, and retention rule reduce the internal investigation needed after an incident. A supplier that provides current provenance and safety contacts avoids repeated emergency searches by each customer. A record that establishes reasonable evaluation, active safeguards, or independent modification can end a claim earlier. The rule converts bespoke forensic reconstruction into standardized retrieval.

Trade secrets require protection rather than categorical exclusion. Rule 26(c)(1)(G) already authorizes orders protecting trade secrets and confidential research, development, or commercial information. Secure expert review, redaction, in camera inspection, access limits, and testing protocols can resolve many disputes without public disclosure. The initial stage seldom requires source code. It asks which version ran, who controlled it, what validation governed, what event occurred, and where the deeper record resides.

Privilege presents a related boundary. Legal advice and attorney work product remain protected. The operator should identify withheld material in a privilege log and produce the underlying nonprivileged facts and business records. A company cannot transform an operational testing program into an inaccessible record by routing every evaluation through counsel, but the precise line follows existing privilege and work-product law. The model statute supplies no new exception.

Privacy and security receive the same design treatment. Requests should use event identifiers, defined windows, pseudonymization, aggregation, and data minimization. Unrelated users' conversations and records fall outside. Security-sensitive controls may be described, tested in a secure environment, or reviewed in camera. These protections can narrow access while preserving the fact that a control, version, evaluation, or incident exists.

The rule therefore rejects two extremes. Public release of every technical artifact would impose unjustified costs and expose protected interests. Complete secrecy would permit a consequential system to define the evidence available against itself. Staged judicial access occupies the administrable middle.

\subsection{Existing Legal Interests Define the Gate}

Henderson emphasizes the heterogeneity of AI activity and argues that some AI-caused losses should receive no tort remedy.\lrfootnote{See \artcite{hendersonRobotLiability2026}.} A disappointed user, an offended reader, a competitor facing lawful innovation, and a patient denied essential care present different interests. Treating every adverse output as compensable would flatten those distinctions.

The front door preserves them at the trigger. A claimant identifies a recognized cause of action or existing equitable entitlement. The verified facts connect the designated deployment and requested record to an element, defense, or remedy. A preference for a different model, a demand for a human conversation, or an assertion that an output felt wrong supplies no information right by itself.

Courts can decide pure legal defenses before expanding production. Jurisdiction, immunity, preemption, a categorical no-duty rule, absence of a private right, limitations apparent from the notice, and failure to identify a protected interest may end the matter. Stage One remains appropriate only to the extent identity or preservation is needed to determine the issue. The statute creates an enforceable right to a first answer for a qualifying notice; it does not create a right to audit any business that uses AI.

The minimum-access principle follows from that gate. A qualifying request receives the information necessary to locate the dispute within an existing legal rule and identify the actors capable of answering it. A claimant invoking defamation connects the request to publication, attribution, fault, or another required issue. An applicant invoking discrimination identifies an adverse process and protected theory. A benefits recipient identifies the entitlement and challenged action. Once those threshold facts place the interest outside governing law, the proceeding ends.

Sectoral designation supplies a second screen. Legislatures can select deployments where affected people predictably lack the facts needed to invoke existing protection. They can omit low-consequence uses, define small-entity procedures, and tailor records to the interest. The resulting statute makes no claim of universal AI exceptionalism.

The access duty is forum neutral and nonwaivable as to the affected person. A vendor agreement can allocate performance and expense between enterprises, but it cannot erase the operator's minimum record or first answer. Confidentiality yields to the defined access right subject to privilege, privacy, trade-secret, and security protection.

Arbitration can select the merits forum when an agreement is otherwise enforceable. The record obligation still arises from deployment and a qualifying request before the private forum needs to compel discovery. That distinction matters because arbitral power over unbound nonparties can be narrower than party discovery. Federal Arbitration Act section 7 authorizes arbitrators to summon a person to attend and bring documents; the Ninth Circuit has held that it does not authorize a stand-alone prehearing document subpoena to a nonparty.\lrfootnote{See \statcite{faaArbitrationSection7}; \casecite[706--09]{cvsVividus2017}.} The statute should therefore require the bound visible operator to maintain the closed local record and contractual retrieval path. An arbitrator can enforce the duty against a bound operator; the design does not depend on creating new compulsory power over an upstream stranger.

Cross-border supply receives the same treatment. A foreign enterprise that operates a covered domestic service remains subject to ordinary jurisdiction, service, comity, and choice-of-law limits. A domestic operator, however, can condition deployment on local retention of the minimum manifest, a reachable supplier agent, material-change notice, and a right to retrieve defined event and provenance fields. The front door solves the information architecture at contracting time instead of implying that Rule 45 reaches every foreign supplier. If a supplier refuses the required interface, the operator can choose another supplier or forgo the designated deployment.

\subsection{Fragmentation Narrows Control Rather Than Eliminating It}

The material-control test is designed for heterogeneous stacks. In a vertically integrated service, one enterprise may own the application, configure the harness, monitor incidents, and issue updates. In an API deployment, the model provider may control base capability and access conditions while the application company controls prompts, retrieval, tools, user population, warnings, and business objectives. In an open-weights deployment, the original publisher may lose continuing visibility and stopping power while a modifier adds the purpose, data, credentials, and interface.

These differences alter who answers after Stage One. They do not defeat the gateway. The visible operator maintains provenance for the system it chose. It identifies upstream components and the rights reserved by each enterprise. The cumulative control test then asks whether a particular lane was risk specific, practically exercisable, and material.

Several recurring defenses follow naturally. A downstream actor that removed safeguards, changed weights, or added dangerous tools may control the relevant path. A user who deliberately defeated a tested constraint may bear comparative or primary responsibility under governing law. A novel cyberattack may interrupt causation, while an unpatched known vulnerability remains part of deployment risk. An upstream provider may show that it could neither observe nor prevent an independent use. A provider that reserved audit, update, suspension, or safety-override rights may have retained a lane. Version and authority records allow the parties to prove these propositions.

Contract labels remain evidence rather than conclusions. An agreement can establish usable rights and information duties. It cannot allocate away a statutory front door owed to an affected person, and the word ``independent'' cannot erase a control the contract actually reserves. Conversely, brand ownership and a general right to terminate a commercial relationship do not establish control over every deployment risk.

The framework leaves room for residual risk. Some harmful outputs will remain unforeseeable after reasonable testing, monitoring, warning, and response. Some interventions will be technically unavailable or would not have changed the event. Ordinary fault and causation standards can produce judgment for the enterprise. A complete deployment record strengthens that result by showing which risks were assessed, which controls were active, what feedback arrived, and why the proposed alternative fails.

\subsection{Product and Speech Classification Remain for the Merits}

AI disputes can implicate design, content, and human expression in different combinations. The front-door rule should not resolve that classification at coverage. A companion application may present product-design questions about age access, memory, engagement objectives, crisis detection, and warnings while a claim imposing liability for an adopted idea raises a distinct First Amendment issue. An employment score may be part of a decision process without becoming a physical product. A generated notice may communicate the operator's decision even when the model has no legal status of its own.

\emph{Garcia} illustrates the separation. At Rule 12, the district court held that the complaint plausibly treated Character.AI as a product because the challenged features concerned design and commercial distribution. It declined on that record to treat the generated output as protected speech where defendants had not identified the relevant expressive choice or speaker.\lrfootnote{See \casecite[1174--80]{garcia2025}.} Those rulings were procedural and claim specific. They neither classified every language model as a product nor denied that human expression mediated through AI can receive protection on another record.

The front door asks for facts that permit the classification inquiry: what conduct the claim challenges; which human or organizational choices the system implemented; what message, if any, an identified person adopted; and whether the requested remedy regulates design, transaction, decision process, or expression. Production remains subject to constitutional limits and ordinary law. If a claim targets protected speech, the court applies the appropriate substantive and procedural safeguards. If it targets a nonexpressive control, the presence of generated words does not automatically remove the design from review.

Anthropomorphic presentation has the same evidentiary status. Calling a system a companion, therapist, or agent cannot create a new legal person. Memory, proactive messages, affection, expertise claims, and relational objectives can nevertheless bear on intended use, foreseeable exposure, warning, and which evaluations are risk specific. The law evaluates the enterprise's own choices without treating the model as speaker, tortfeasor, or liability sink by default.

The boundary completes the proposal. The front door is classification neutral but fact enabling. It creates no preference for liability and no immunity through metaphor. It permits courts to choose the correct substantive category on a record that reveals the deployment rather than on a label selected before the relevant facts were available.
